\subsection{Hilbert--Schmidt integral operators}

In this issue, X and Y are locally compact topological spaces. Let \(\mu\) be a positive measure on X and \(\nu\) be a positive measure on Y. We denote by \(\mathrm{L}^{2}(\mathrm{X})\) , \(\mathrm{L}^{2}(\mathrm{Y})\) and \(\mathrm{L}^{2}(\mathrm{X} \times \mathrm{Y})\) the spaces \(\mathrm{L}^{2}(\mathrm{X}, \mu)\) , \(\mathrm{L}^{2}(\mathrm{Y}, \nu)\) and \(\mathrm{L}^{2}(\mathrm{X} \times \mathrm{Y}, \mu \otimes \nu)\) respectively, and likewise for \(\mathcal{L}^{2}(\mathrm{X})\) , \(\mathcal{L}^{2}(\mathrm{Y})\) and \(\mathcal{L}^{2}(\mathrm{X} \times \mathrm{Y})\) .

We recall (cf. No. 4 of III, p. 26) that for every \(N \in \mathcal{L}^{2}(X \times Y)\) there exists a unique continuous linear map \(u_{\mathrm{N}}\) from \(L^{2}(X)\) into \(L^{2}(Y)\) such that

\[
\langle g \mid u _ {\mathrm{N}} (f) \rangle = \int_ {\mathrm{X} \times \mathrm{Y}} \overline {{g (y)}} \mathrm{N} (x, y) f (x) d (\mu \otimes \nu) (x, y)\tag{12}
\]

for all \(f \in \mathrm{L}^{2}(\mathrm{X})\) and every \(g \in \mathrm{L}^{2}(\mathrm{Y})\) . The map \(u_{N}\) is compact (cor. 1 of III, p. 33). The map \(N \mapsto u_{N}\) induces, by passage to quotients, a continuous and injective linear map from \(\mathrm{L}^{2}(\mathrm{X} \times \mathrm{Y})\) into \(\mathcal{L}(\mathrm{L}^{2}(\mathrm{X});\mathrm{L}^{2}(\mathrm{Y}))\) (prop. 5 of III, p. 30).

We denote by \(\theta\) the unique linear map from \(\mathcal{L}^2 (\mathrm{X})\otimes \mathcal{L}^2 (\mathrm{Y})\) into \(\mathcal{L}^2 (\mathrm{X}\times \mathrm{Y})\) which, for all \(f\in \mathcal{L}^2 (\mathrm{X})\) and \(g\in \mathcal{L}^2 (\mathrm{Y})\), sends \(f\otimes g\) to the function defined by \((x,y)\mapsto f(x)g(y)\) .

Lemma 10. --- The map \(\theta\) defines by passage to quotients a linear map \(\widetilde{\theta}\) from \(\mathrm{L}^2 (\mathrm{X})\otimes \mathrm{L}^2 (\mathrm{Y})\) into \(\mathrm{L}^2 (\mathrm{X}\times \mathrm{Y})\) and there exists a unique isometric isomorphism from \(\mathrm{L}^2 (\mathrm{X})\widehat{\otimes}_2\mathrm{L}^2 (\mathrm{Y})\) onto \(\mathrm{L}^2 (\mathrm{X}\times \mathrm{Y})\) which coincides with it on \(\mathrm{L}^2 (\mathrm{X})\otimes \mathrm{L}^2 (\mathrm{Y})\) .

The first assertion is elementary. Let us prove the second.

Let \((f_i)_{i\in \mathrm{I}}\) and \((g_j)_{j\in \mathrm{J}}\) be orthonormal bases of \(\mathrm{L}^2 (\mathrm{X})\) and \(\mathrm{L}^2 (\mathrm{Y})\) respectively. The family \((\overline{f}_i\otimes g_j)_{(i,j)\in \mathrm{I}\times \mathrm{J}}\) is an orthonormal basis of \(\mathrm{L}^2 (\mathrm{X})\widehat{\otimes}_2\mathrm{L}^2 (\mathrm{Y})\) (EVT, V, p. 29, cor. 1) and the family \((\widetilde{\theta} (\overline{f}_i\otimes g_j))_{(i,j)\in \mathrm{I}\times \mathrm{J}}\) is orthonormal in \(\mathrm{L}^2 (\mathrm{X}\times \mathrm{Y})\). The map \(\widetilde{\theta}\) extends therefore by continuity to an isometric linear map from \(\mathrm{L}^2 (\mathrm{X})\widehat{\otimes}_2\mathrm{L}^2 (\mathrm{Y})\) into \(\mathrm{L}^2 (\mathrm{X}\times \mathrm{Y})\). It remains to show that this extension is surjective.

For this, it suffices to prove that the image of \(\widetilde{\theta}\) is dense in \(\mathrm{L}^{2}(\mathrm{X}\times\mathrm{Y})\). Let N be an element of \(\mathrm{L}^{2}(\mathrm{X}\times\mathrm{Y})\) orthogonal to the image of \(\widetilde{\theta}\) . For all \(i\in I\) and \(j\in J\) , we have \(\langle g_{j}\mid u_{\mathrm{N}}(f_{i})\rangle=\langle\widetilde{\theta}(\overline{f}_{i}\otimes g_{j})\mid\mathrm{N}\rangle=0\). From formula (12), it follows that \(u_{N}=0\) and therefore N=0.

The preceding lemma establishes the assertion stated in EVT, V, p. 29, example 2.

In what follows, we will identify \(\mathrm{L}^{2}(\mathrm{X})\widehat{\otimes}_{2}\mathrm{L}^{2}(\mathrm{Y})\) and \(\mathrm{L}^{2}(\mathrm{X}\times\mathrm{Y})\) by the isometric isomorphism of Lemma 10.

PROPOSITION 19. --- The map \(N\mapsto u_{\mathrm{N}}\) is an isometric isomorphism from \(\mathrm{L}^{2}(\mathrm{X}\times\mathrm{Y})\) onto the space \(\mathcal{L}_{2}(\mathrm{L}^{2}(\mathrm{X});\mathrm{L}^{2}(\mathrm{Y}))\) of Hilbert--Schmidt maps from \(\mathrm{L}^{2}(\mathrm{X})\) into \(\mathrm{L}^{2}(\mathrm{Y})\) .

The linear map from \(\mathrm{L}^{2}(\mathrm{Y})\otimes\overline{\mathrm{L}^{2}(\mathrm{X})}\) into \(\mathcal{L}_{2}(\mathrm{L}^{2}(\mathrm{X});\mathrm{L}^{2}(\mathrm{Y}))\) that sends \(g\otimes f\) to the Hilbert--Schmidt map \(h\mapsto\langle f|h\rangle g\) extends to an isometric isomorphism \(\theta_{1}\) of \(\mathrm{L}^{2}(\mathrm{Y})\widehat{\otimes}_{2}\overline{\mathrm{L}^{2}(\mathrm{X})}\) onto \(\mathcal{L}_{2}(\mathrm{L}^{2}(\mathrm{X});\mathrm{L}^{2}(\mathrm{Y}))\) (EVT, V, p. 52, th. 1).

Additionally, let us denote by \(\theta_{2}\) the isometric isomorphism of \(\mathrm{L}^{2}(\mathrm{X})\widehat{\otimes}_{2}\mathrm{L}^{2}(\mathrm{Y})\) onto \(\mathrm{L}^{2}(\mathrm{Y})\widehat{\otimes}_{2}\overline{\mathrm{L}^{2}(\mathrm{X})}\) which sends \(f\otimes g\) to \(g\otimes \overline{f}\) for all \(f\in \mathrm{L}^2 (\mathrm{X})\) and every \(g\in \mathrm{L}^2 (\mathrm{Y})\) .

The linear map \(\theta_{3} = \theta_{1} \circ \theta_{2}\) is identified with an isometric isomorphism of \(\mathrm{L}^{2}(\mathrm{X} \times \mathrm{Y})\) onto \(\mathcal{L}_{2}(\mathrm{L}^{2}(\mathrm{X}); \mathrm{L}^{2}(\mathrm{Y}))\) .

Let \(f \in \mathrm{L}^2(\mathrm{X})\) and \(g \in \mathrm{L}^2(\mathrm{Y})\) and let \(\mathrm{N}\) be the element of \(\mathrm{L}^2(\mathrm{X} \times \mathrm{Y})\) identified with \(f \otimes g\). By INT, V, p. 95, § 8, no. 3, cor. 2 and formula (12), we have

\[
\langle g _ {1} \mid u _ {\mathrm{N}} (f _ {1}) \rangle = \langle \overline {{f}} \mid f _ {1} \rangle \langle g _ {1} \mid g \rangle
\]

for all \(f_1 \in \mathrm{L}^2(\mathrm{X})\) and \(g_1 \in \mathrm{L}^2(\mathrm{Y})\) . The map \(u = \theta_3(\mathrm{N})\) is the linear map \(\theta_1(g \otimes \overline{f})\); it therefore satisfies \(u(h) = \langle \overline{f} | h \rangle g\) for all \(h \in \mathrm{L}^2(\mathrm{X})\). Consequently, for all \(f_1 \in \mathrm{L}^2(\mathrm{X})\) and \(g_1 \in \mathrm{L}^2(\mathrm{Y})\) , it follows that

\[
\langle g _ {1} \mid u (h _ {1}) \rangle = \langle \overline {{f}} \mid h _ {1} \rangle \langle g _ {1} \mid g \rangle = \langle g _ {1} \mid u _ {\mathrm{N}} (h _ {1}) \rangle ,
\]

whence \(\theta_3(N) = u_N\) . Since \(\theta_3\) and \(N \mapsto u_N\) are continuous, we conclude that \(\theta_3(N) = u_N\) for all \(N \in L^2(X \times Y)\) , which concludes the proof.

COROLLARY. --- For every \(N \in L^{2}(X \times Y)\) , the linear map \(u_{N}\) is a Hilbert--Schmidt operator and we have \(\mathrm{Tr}(u_{\mathrm{N}}^{*}u_{\mathrm{N}}) = \|N\|^{2}\) .

Indeed, one has \(\| u\| _2 = \sqrt{\mathrm{Tr}(u^*u)}\) for all \(u\in \mathcal{L}_2(\mathrm{L}^2 (\mathrm{X});\mathrm{L}^2 (\mathrm{Y}))\) .

Note. --- Let \(N \in L^{2}(X \times Y)\) . By Corollary 2 of IV, p. 150, there exist a countable set I, orthonormal families \((f_{i})_{i \in I}\) in \(L^{2}(X)\) and \((g_{i})_{i \in I}\) in \(L^{2}(Y)\) , and a family \((\alpha_{i})_{i \in I}\) in \(\mathbf{R}_{+}^{*}\) such that

\[
u _ {\mathrm{N}} (f) = \sum_ {i \in \mathrm{I}} \alpha_ {i} \langle f _ {i} | f \rangle g _ {i}
\]

for all \(f \in \mathrm{L}^2(\mathrm{X})\), where the series converges in \(\mathrm{L}^2(\mathrm{Y})\) . By cor. 4 of IV, p. 165 and the corollary above, we have

\[
\sum_ {i \in \mathrm{I}} \alpha_ {i} ^ {2} = \mathrm{Tr} (u _ {\mathrm{N}} ^ {*} u _ {\mathrm{N}}) = \| u _ {\mathrm{N}} \| _ {2} ^ {2} = \int_ {\mathrm{X} \times \mathrm{Y}} | \mathrm{N} (x, y) | ^ {2} d (\mu \otimes \nu) (x, y).
\]

Additionally, let us note that \(h_{i,j} = \overline{f}_{i} \otimes g_{j} \in \mathrm{L}^{2}(\mathrm{X} \times \mathrm{Y})\) for all \((i, j) \in \mathrm{I} \times \mathrm{I}\) . We then have

\[
\mathrm{N} = \sum_ {i \in \mathrm{I}} \alpha_ {i} h _ {i, i}
\]

in \(\mathrm{L}^2 (\mathrm{X}\times \mathrm{Y})\) .

Indeed, let \((f_{j})_{j\in\mathrm{J}}\) and \((g_{k})_{k\in\mathrm{K}}\) be orthonormal bases of \(\mathrm{L}^{2}(\mathrm{X})\) and \(\mathrm{L}^{2}(\mathrm{Y})\) , respectively, extending the families \((f_{i})_{i\in\mathrm{I}}\) and \((g_{i})_{i\in\mathrm{I}}\) . Set \(h_{j,k}=\overline{f}_{j}\otimes g_{k}\) for all \((j,k)\in\mathrm{J}\times\mathrm{K}\) . By Lemma 10, the family \((h_{j,k})_{(j,k)\in\mathrm{J}\times\mathrm{K}}\) is an orthonormal basis of \(\mathrm{L}^{2}(\mathrm{X}\times\mathrm{Y})\) . We have \(\langle h_{j,k}\mid\mathrm{N}\rangle=\langle g_{k}\mid u_{\mathrm{N}}(f_{j})\rangle\) for all \((j,k)\in\mathrm{J}\times\mathrm{K}\) . If \(j\notin I\) this quantity is zero. If \(j\in I\) , it is equal to \(\alpha_{j}\langle g_{k}\mid g_{j}\rangle\) , hence is zero unless k=j, in which case \(\langle h_{j,j}\mid N\rangle=\alpha_{j}\) . Consequently

\[
\mathrm{N} = \sum_ {(j, k) \in \mathrm{J} \times \mathrm{K}} \left\langle h _ {j, k} \mid \mathrm{N} \right\rangle h _ {j, k} = \sum_ {i \in \mathrm{I}} \alpha_ {i} h _ {i, i}.
\]

